% ============================================================================
%  chapter3.tex
%  Building Abstractions with Structs
% ============================================================================

\chapter{Building Abstractions with Structs}

\section{From Scalars to Records}

Chapter~1 gave us bytes, addresses, and pointers, and Chapter~2
opened the registers that hold them while the machine runs. With only
scalar
types (\code{int}, \code{char}, \code{double}) we can already do real
work, but real systems are full of \emph{aggregates}: a connection has
a socket, a buffer, a state, a timeout; a record has a key and a
value; a node has data and links to its neighbors. In C, the
\emph{struct} is the tool that groups related data into one object
whose bytes live contiguously in memory.

A struct is a \emph{layout contract}: it tells the compiler exactly how
to arrange a set of fields one after another in memory, so that the
whole group can be passed around, pointed at, and stored as a single
unit. Everything else in this book---linked lists, hash tables, trees,
graphs, sockets---is, at bottom, structs pointing at other structs.

\section{Declaring and Using a Struct}

\begin{lstlisting}[caption={A minimal struct: a 2D point.}]
#include <stdio.h>

struct point {
    int x;
    int y;
};

int main(void) {
    struct point origin;       /* the struct is a variable like any other */

    origin.x = 0;              /* fields are accessed with the dot */
    origin.y = 0;

    printf("origin = (%d, %d)\n", origin.x, origin.y);
    return 0;
}
\end{lstlisting}

The declaration \code{struct point \{ int x; int y; \};} introduces a
new type named \code{struct point}. It does \emph{not} allocate any
memory; it merely describes the shape. Only when we declare a variable
of that type---\code{struct point origin;}---does the compiler reserve
the bytes. The memory reserved is exactly the sum of the field sizes,
plus whatever padding alignment requires, as we are about to see.

\subsection{typedef: Giving the Struct a Short Name}

Writing \code{struct point} everywhere is verbose, and in this book we
will write a great many struct types. \code{typedef} introduces an
alias for an existing type, so we can say \code{point} and mean
\code{struct point}. From here on we will declare our structs with a
\code{typedef} and refer to them by the bare name.

\begin{lstlisting}[caption={A typedef names the struct; a companion pointer typedef follows.}]
#include <stdio.h>

typedef struct point {
    int x;
    int y;
} point;

typedef struct node {
    int          value;
    struct node *next;   /* inside the definition we still use "struct node" */
} node;

int main(void) {
    point origin;        /* "point" is an alias for "struct point" */

    origin.x = 0;
    origin.y = 0;

    printf("(%d, %d)\n", origin.x, origin.y);
    return 0;
}
\end{lstlisting}

Two notes on discipline. First, the tag name \code{point} in
\code{typedef struct point \{...\} point;} is optional but useful for
self-documentation. Second, \code{typedef} does \emph{not} create a
new type; it creates an alias. The rules of argument passing, copying,
and layout apply to the alias exactly as they do to the original. And a
pointer typedef, \code{typedef struct node *node\_ptr;}, is a style
choice many codebases make; we avoid it in this course, because
hiding a pointer inside a typedef makes it easy to forget whether you
hold a pointer or a value. Pointers stay explicit.

\section{Memory Layout of a Struct}

Because a struct's fields are laid out in declaration order, we can
draw the layout of \code{struct point} with perfect precision:

\begin{verbatim}
   struct point { int x; int y; }   -- size 8 bytes

   offset:  0     1     2     3     4     5     6     7
           +-----+-----+-----+-----+-----+-----+-----+-----+
           |           x           |           y           |
           +-----+-----+-----+-----+-----+-----+-----+-----+
            ^                       ^
            address of origin       address of origin.y
            == address of origin.x
\end{verbatim}

The address of a struct is the address of its first field. The two are
the same number, which is exactly why \code{\&origin.x} and
\code{\&origin} compare equal. Every later field sits at a fixed
\emph{offset} from the start of the struct.

\begin{lstlisting}[caption={A struct is an int and another int, one after another.}]
#include <stdio.h>

struct point {
    int x;
    int y;
};

int main(void) {
    struct point origin;

    origin.x = 10;
    origin.y = 20;

    printf("sizeof(struct point) = %zu\n", sizeof(struct point));
    printf("&origin      = %p\n", (void *) &origin);
    printf("&origin.x    = %p\n", (void *) &origin.x);
    printf("&origin.y    = %p\n", (void *) &origin.y);
    printf("offset of y  = %ld\n",
           (long) ((char *) &origin.y - (char *) &origin.x));

    return 0;
}
\end{lstlisting}

We printed the offset by subtracting two \code{char*} pointers. Pointer
subtraction between two pointers into the same array of \code{char}
yields the number of \code{char} elements between them, which is the
offset in bytes. This is a legitimate, explicit technique; note that we
cast to \code{char*} precisely so that the arithmetic counts bytes and
not elements.

\section{Data Alignment: Why the CPU Demands It}

Here is the fact about modern hardware that explains everything that
looks ``wasteful'' about structs: a CPU reads memory in aligned units.
On x86-64, a 4-byte \code{int} is best placed at an address divisible
by 4; an 8-byte \code{double} at an address divisible by 8. When a
multi-byte object straddles a boundary, the hardware may need two
fetches instead of one, or---on stricter architectures such as some
RISC designs---may raise a fault outright.

\begin{verbatim}
   Correctly aligned (good):

   address: 1000  1001  1002  1003  1004  1005  1006  1007
           +-----+-----+-----+-----+-----+-----+-----+-----+
           |           an int at 1000         |  (next obj) |
           +-----+-----+-----+-----+-----+-----+-----+-----+
            ^ int starts at a multiple of 4

   Misaligned (bad):

   address: 1002  1003  1004  1005  1006  1007  1008  1009
           +-----+-----+-----+-----+-----+-----+-----+-----+
           | (prev) |      an int starting at 1002      |   |
           +-----+-----+-----+-----+-----+-----+-----+-----+
            ^ 1002 is not a multiple of 4: the int straddles a
              boundary and may need two fetches to be read.
\end{verbatim}

To honor alignment, the compiler may insert \emph{padding} bytes
between struct fields. These bytes exist, are addressable, and are
completely invisible to you unless you look at the layout directly.
They cost memory but buy speed and, on many machines, correctness.

\begin{ruleofthumb}[Never assume a field's offset]
The C standard guarantees fields are laid out in declaration order, but
\emph{not} their exact offsets, because padding depends on the machine
and the compiler. Any code that assumes \code{int} sits at offset 0, or
that a struct is exactly the sum of its parts, is broken by definition.
Use \code{sizeof} and \code{offsetof} if you must know the numbers.
\end{ruleofthumb}

\section{Padding: The Hidden Bytes}

Consider a structure whose fields have different alignments:

\begin{lstlisting}[caption={A struct with mixed field sizes.}]
struct entry {
    char   tag;     /* 1 byte,  alignment 1 */
    int    count;   /* 4 bytes, alignment 4 */
    short  level;   /* 2 bytes, alignment 2 */
};
\end{lstlisting}

The compiler must place each field at an offset divisible by its own
alignment. The resulting layout, with padding marked \code{//}:

\begin{verbatim}
   struct entry -- size 12, with 5 padding bytes hidden inside

   offset:  0     1     2     3     4     5     6     7
           +-----+-----+-----+-----+-----+-----+-----+-----+
           | tag | //  | //  | //  |        count         |
           +-----+-----+-----+-----+-----+-----+-----+-----+
   offset:  8     9    10    11
           +-----+-----+-----+-----+
           | level |     //    | // |
           +-----+-----+-----+-----+

   - tag    at offset 0  (alignment 1: any address)
   - count  at offset 4  (must be a multiple of 4; bytes 1-3 are padding)
   - level  at offset 8  (must be a multiple of 2)
   - total raw payload:  1 + 4 + 2 = 7 bytes
   - struct size: 12     (rounded up to a multiple of 4, the largest
                           alignment in the struct, so that arrays of
                           these structs stay aligned)
\end{verbatim}

The struct ended up at 12 bytes for 7 bytes of real data. Five of the
twelve bytes are silent. This is not a bug; it is the machine telling
you how it wants to be spoken to. But notice the waste is not forced:
nothing about the fields themselves demands that layout. If we
redeclare the fields in the right order, the compiler can satisfy every
alignment without a single padding byte between fields.

\section{Field Ordering to Shrink Structs}

Place fields in decreasing order of their alignment. The largest
alignment dominates; once the big fields are placed, the small ones
slot into the gaps.

\begin{lstlisting}[caption={The same data, reordered.}]
struct entry {
    int    count;   /* 4 bytes, alignment 4 */
    short  level;   /* 2 bytes, alignment 2 */
    char   tag;     /* 1 byte,  alignment 1 */
};
\end{lstlisting}

\begin{verbatim}
   struct entry, reordered -- size 8

   offset:  0     1     2     3     4     5     6     7
           +-----+-----+-----+-----+-----+-----+-----+-----+
           |        count         | level | tag | //  |
           +-----+-----+-----+-----+-----+-----+-----+-----+

   - count  at offset 0
   - level  at offset 4
   - tag    at offset 6
   - one padding byte at offset 7 (to round the size up to 8)
   - raw payload still 7 bytes, but the struct is now 8, not 12.
\end{verbatim}

The same information shrank from 12 bytes to 8 bytes, a third smaller,
for free. In a table of a million entries that is 4~MB saved, and on a
cache-bound hot path it is substantially more than a memory saving.
Field ordering is the cheapest optimization in C and one that beginners
never think about.

\begin{insight}[The struct is a memory-layout contract, not just a type]
Reading a struct definition is reading a memory diagram. Before you
write a struct for a hot path or a wire format, order its fields from
largest alignment to smallest, and get into the habit of sketching the
layout by hand. The machine will reward you.
\end{insight}

\section{offsetof and sizeof}

C gives us two operators for interrogating layouts at compile time.
\code{sizeof(type)} returns the total size including trailing padding;
\code{offsetof(type, field)} returns the byte offset of a field.

\begin{lstlisting}[caption={Reading the layout without guessing.}]
#include <stdio.h>
#include <stddef.h>

struct entry {
    int    count;
    short  level;
    char   tag;
};

int main(void) {
    printf("sizeof(struct entry)      = %zu\n", sizeof(struct entry));
    printf("offsetof count            = %zu\n", offsetof(struct entry, count));
    printf("offsetof level            = %zu\n", offsetof(struct entry, level));
    printf("offsetof tag              = %zu\n", offsetof(struct entry, tag));
    return 0;
}
\end{lstlisting}

\code{offsetof} is what lets a generic serializer walk a struct without
special-casing each field: given a pointer to the struct, the address
of a field is simply
\code{(char *) base + offsetof(struct entry, level)}. We will use
exactly this idea when we build a data serialization layer in the
chapters on file I/O and networking.

\begin{ruleofthumb}[Arrays of structs inherit the alignment]
The compiler rounds a struct's size up to a multiple of its strictest
member so that an array of such structs has every element aligned. That
trailing padding is why \code{sizeof} can exceed the sum of the
\code{offsetof} values plus the last field's size. Treat
\code{sizeof(struct)} as authoritative; never compute it by hand.
\end{ruleofthumb}

\section{Structs and Pointers: The Arrow Operator}

Pointers to structs are everywhere, because structs are the units we
link together. To reach a field through a pointer, C offers the
\emph{arrow operator}: \code{p->field} means \code{(*p).field}.

\begin{lstlisting}[caption={Accessing fields through a pointer.}]
#include <stdio.h>

struct point {
    int x;
    int y;
};

/* Move the point that p points at by (dx, dy). */
void translate(struct point *p, int dx, int dy) {
    p->x = p->x + dx;   /* equivalent to (*p).x = (*p).x + dx; */
    p->y = p->y + dy;
}

int main(void) {
    struct point here;
    struct point *cursor;

    here.x = 3;
    here.y = 4;

    cursor = &here;               /* cursor now points at here */
    translate(cursor, 5, 6);      /* here becomes (8, 10)     */

    printf("here = (%d, %d)\n", here.x, here.y);
    printf("via pointer = (%d, %d)\n", cursor->x, cursor->y);
    return 0;
}
\end{lstlisting}

\begin{verbatim}
   Memory state after cursor = &here:

        +--------------------+
        |      here          |
        |  x = 3  y = 4      |
        +--------------------+
            ^
            |
        cursor  (a struct point* holding &here)
\end{verbatim}

Why pointers instead of passing the struct by value? A pointer passes 8
bytes no matter how large the struct is; copying a big struct by value
copies every byte. More importantly, a function that receives a pointer
can \emph{modify} the caller's struct, exactly as \code{translate} did
above. If you want a function to change a struct, pass its address.

\section{const Correctness with Structs}

\code{const} is the compiler's way of recording your intent, and in
struct-heavy code it earns its keep. The keyword has two different
places to attach: to the pointed-to struct (\code{const struct person
*p}, read-only object) or to the pointer itself (\code{struct person
*const p}, a pointer that cannot be reassigned). The first is the one
we use constantly in function signatures.

\begin{lstlisting}[caption={const expresses ``read-only'' on function boundaries.}]
#include <stdio.h>

typedef struct point {
    int x;
    int y;
} point;

/* "p" points at a struct the caller promises not to modify through us. */
int distance_squared(const point *p) {
    return p->x * p->x + p->y * p->y;
}

/* The pointer itself is const: inside this function, q cannot be reassigned. */
void print_point(point *const q) {
    printf("(%d, %d)\n", q->x, q->y);
}

int main(void) {
    point origin = {0, 0};

    print_point(&origin);
    printf("dist^2 = %d\n", distance_squared(&origin));
    return 0;
}
\end{lstlisting}

Read the signatures as promises. \code{distance\_squared} promises it
will not modify the point; a caller can therefore hand it any point,
confident nothing is corrupted. \code{const} does not make memory
physically read-only---it is a compile-time contract, and a
\code{const} pointer to a struct is still an ordinary 8-byte address.
But the contract is enforced: writing through a \code{const} pointer is
a compile error, which turns a whole class of lifetime and aliasing
bugs into loud, immediate failures at build time.

\begin{ruleofthumb}[Declare read-only parameters const]
Any function that receives a pointer to a struct and does not modify
it should declare that parameter \code{const}. It costs nothing, it
documents the interface, and it lets the compiler catch accidental
writes. In Chapter~4 every ``print'' and ``lookup'' function takes
\code{const} pointers; every ``insert'' and ``remove'' does not.
\end{ruleofthumb}

\section{Structs on the Heap}

So far our structs have lived on the stack. But the stack frame dies
with the function, so any struct that must outlive its creator, or
whose count is known only at runtime, is allocated on the heap with
\code{malloc}. This is the moment when structs and pointers become
genuinely powerful: \code{malloc} returns the address of a fresh,
uninitialized region big enough for the struct, and we then fill its
fields through that address.

\begin{lstlisting}[caption={A struct living on the heap, with ownership handled.}]
#include <stdio.h>
#include <stdlib.h>

struct account {
    int   id;
    long  balance;
};

struct account *account_new(int id) {
    struct account *a;

    a = malloc(sizeof(struct account));   /* ask for exactly one struct */
    if (a == NULL) {
        return NULL;
    }

    a->id = id;
    a->balance = 0;

    return a;
}

int main(void) {
    struct account *acc;

    acc = account_new(7);
    if (acc == NULL) {
        fprintf(stderr, "allocation failed\n");
        exit(1);
    }

    printf("account %d has balance %ld\n", acc->id, acc->balance);

    free(acc);       /* we own it, so we free it; no more acc->id after this */
    return 0;
}
\end{lstlisting}

Three details deserve emphasis, because they are the core discipline of
every data structure we build in this book:

\begin{itemize}
  \item \code{malloc(sizeof(struct account))} asks for exactly the
    struct's size---padding included---so the layout contract carries
    over to the heap without us ever computing byte counts.
  \item The \code{if (a == NULL)} check is the contract of every
    allocation. Heap exhaustion is a real, reportable outcome.
  \item Ownership is explicit: \code{account\_new} allocates, and the
    caller frees. If \code{acc} is later embedded in a linked list or a
    hash table, the data structure becomes the owner, and it is the
    data structure's \emph{destroy} function that frees it.
\end{itemize}

\section{Nested Structs}

A struct may contain another struct \emph{by value}. The inner struct's
bytes are laid out inline, within the outer struct, at the offset its
alignment requires. The memory diagram makes the word ``nested'' literal:

\begin{lstlisting}[caption={A struct containing a struct.}]
#include <stdio.h>

typedef struct point {
    int x;
    int y;
} point;

typedef struct rect {
    point  top_left;     /* a whole point, embedded */
    point  bottom_right; /* another whole point, embedded */
} rect;
\end{lstlisting}

\begin{verbatim}
   struct rect -- size 16

   offset:  0     1     2     3     4     5     6     7
           +-----+-----+-----+-----+-----+-----+-----+-----+
           |      top_left.x       |      top_left.y       |
           +-----+-----+-----+-----+-----+-----+-----+-----+
   offset:  8     9    10    11    12    13    14    15
           +-----+-----+-----+-----+-----+-----+-----+-----+
           |   bottom_right.x     |   bottom_right.y       |
           +-----+-----+-----+-----+-----+-----+-----+-----+

   top_left   occupies bytes 0..7
   bottom_right occupies bytes 8..15
\end{verbatim}

Access uses the dot twice: \code{box.top\_left.x}. There is no arrow
involved unless you reach the nested field through a pointer to the
outer struct, in which case it is \code{box->top\_left.y}. The nested
struct has no separate identity; it is a set of fields at fixed offsets
inside its parent. This is how we build hierarchy from flat memory.

\section{Flexible Array Members: The Variable-Length Tail}

C99 introduced the \emph{flexible array member}: a struct whose last
field is an unsized array \code{data[]}. The struct itself occupies
only the fixed fields; the array bytes are tacked on at the end by
whichever allocation you make. This is exactly how real kernels and
libraries carry a length plus a payload in one allocation, and it is
the pattern behind messages, packets, and buffers throughout this book.

\begin{lstlisting}[caption={One allocation holds header and payload together.}]
#include <stdio.h>
#include <stdlib.h>
#include <string.h>

typedef struct message {
    unsigned int len;      /* number of valid bytes in data */
    char data[];           /* flexible array member: no size here */
} message;

int main(void) {
    message *msg;
    size_t payload;

    payload = 5;                        /* we want to carry 5 bytes */
    msg = malloc(sizeof(message) + payload);
    if (msg == NULL) {
        exit(1);
    }

    msg->len = payload;
    memcpy(msg->data, "hello", payload);

    printf("message: len=%u data=%.*s\n", msg->len, (int) msg->len, msg->data);

    free(msg);                          /* one allocation, one free */
    return 0;
}
\end{lstlisting}

\begin{verbatim}
   The allocation from malloc(sizeof(message) + 5):

   +----------------+------------------------------+
   |   len  (4B)    |   data[0..4]  ("hello")      |
   +----------------+------------------------------+
   ^  sizeof(message) = 4            ^  exactly payload bytes

   The whole thing is ONE contiguous allocation, freed with one free().
   No pointer chase, no leak, no fragmentation from a second malloc.
\end{verbatim}

The rules are strict: the flexible array must be the last member, there
must be at least one named member before it, and you cannot
\emph{store} an array of these structs---you must allocate each one
individually, because its size is not fixed at compile time. Used
honestly, it removes an entire class of double-allocation and
free-order bugs from protocols, virtual machines, and caches.

\section{Self-Referential Structures}

A struct may contain a pointer to an object of its own type. The type
is incomplete during its own definition, but a \emph{pointer} to it is
always allowed, because a pointer is a fixed-size address regardless of
what it targets. This one trick---structs pointing at themselves---is
the seed from which linked lists, trees, and graphs all grow.

\begin{lstlisting}[caption={The birth of every linked structure.}]
struct node {
    int value;
    struct node *next;   /* pointer to another node, or NULL */
};
\end{lstlisting}

\begin{verbatim}
   Three nodes linked into a chain:

   head
    |
    v
   +---------+      +---------+      +---------+
   | value=10|      | value=20|      | value=30|
   | next ---|----->| next ---|----->| next ---|-----> NULL
   +---------+      +---------+      +---------+
      ^                 ^                 ^
      node A            node B            node C
      (heap)            (heap)            (heap)

   Each box is a struct node. The last "next" points at NULL, marking
   the end of the chain. The whole list is named by the single pointer
   "head".
\end{verbatim}

The list is a web of arrows, exactly the diagrams we drew in
Chapter~1. A linked list is not a special language feature; it is a
pointer stored inside a struct pointing at another instance of that
struct. That realization, once felt, unlocks every container in the
standard toolbox.

\section{Copying and Returning Structs}

Unlike arrays, structs can be copied and returned by value. The
compiler copies the whole layout, padding included. A struct returned
by value is a complete, independent object.

\begin{lstlisting}[caption={Structs copy and return by value; arrays do not.}]
#include <stdio.h>

struct point {
    int x;
    int y;
};

struct point make_point(int x, int y) {
    struct point p;

    p.x = x;
    p.y = y;
    return p;          /* a copy of p is handed back; safe, p is local */
}

int main(void) {
    struct point a = make_point(2, 3);
    struct point b;

    b = a;             /* b is now an independent copy of a */
    b.x = 100;

    printf("a = (%d, %d)\n", a.x, a.y);   /* still (2, 3) */
    printf("b = (%d, %d)\n", b.x, b.y);   /* (100, 3)    */
    return 0;
}
\end{lstlisting}

Returning a \emph{struct} by value is safe; returning a pointer to a
local struct is the dangling-pointer bug of Chapter~1. The rule of
thumb: small structs (a few words) are fine by value; large ones are
cheaper as pointers, but then you own the lifetime problem that comes
with the heap.

\begin{insight}[Shallow vs.\ deep copy]
\code{b = a} copies the struct's bytes verbatim. If the struct contains
a pointer (like our \code{node} above), the \emph{pointer} is copied,
not the node it points at: \code{b.next} and \code{a.next} name the
same memory. That is a shallow copy, and it is usually what we want in
containers. Deep copying---duplicating everything reachable through
pointers---is a deliberate, manual operation we will write by hand when
we build graphs in Chapter~4.
\end{insight}

\section{Array of Structs vs.\ Struct of Arrays}

When you hold many records, there are two ways to arrange them in
memory, and the choice is a measurable performance decision that
systems programmers make on every hot path.

\textbf{Array of Structs} (AoS) interleaves the fields: element 0's all
fields, then element 1's all fields. \textbf{Struct of Arrays} (SoA)
groups like fields together: all \code{x} values, then all \code{y}
values.

\begin{verbatim}
   Array of Structs (AoS):              Struct of Arrays (SoA):

   point points[4];                     struct points {
                                            int xs[4];
       [x0 y0 x1 y1 x2 y2 x3 y3]            int ys[4];
                                        };

   each element is contiguous:          all xs live together:
   points[i].x and points[i].y         points.xs[i]
   are adjacent bytes                  points.ys[i] are far apart

   good when you touch both x and y     good when a pass touches only x
   together (all fields of a record)    (a single field across records)
\end{verbatim}

The modern memory hierarchy makes this matter. The CPU loads memory in
cache lines (typically 64 bytes). AoS is ideal when each record is used
wholesale---the whole record arrives in one line. SoA is ideal when an
algorithm scans one field of every record, because consecutive accesses
land in consecutive addresses, and the fields you never touch are never
pulled into cache. Scientific code, game engines, and database column
stores all exploit SoA.

\begin{insight}[Which layout does your access pattern want?]
Ask: do my loops walk \emph{records} or \emph{fields}? Walking records
wants AoS; walking one field at a time wants SoA. There is no
universal winner, and getting this right is often worth more than any
micro-optimization inside a loop.
\end{insight}

\section{Bit Fields: Squeezing Flags into Bytes}

For control registers, file formats, and network headers, C offers
\emph{bit fields}: members whose width is specified in bits. The layout
of bit fields is implementation-defined in painful detail, so we treat
them as a specialized tool, not a general one. They are invaluable
when a hardware or protocol spec forces bits to be packed.

\begin{lstlisting}[caption={Bit fields: several flags in one storage unit.}]
#include <stdio.h>

struct flags {
    unsigned int valid  : 1;   /* one bit  */
    unsigned int dirty  : 1;   /* one bit  */
    unsigned int kind   : 3;   /* three bits (0..7) */
};

int main(void) {
    struct flags f;

    f.valid = 1;
    f.dirty = 0;
    f.kind  = 5;

    printf("valid=%u dirty=%u kind=%u\n", f.valid, f.dirty, f.kind);
    return 0;
}
\end{lstlisting}

Notice we declared them \code{unsigned int}. A signed one-bit field can
only hold $-1$ or $0$, which is almost never what anyone means.

\begin{ruleofthumb}[Bit fields are for hardware and wire formats]
If your struct is destined for a control register or a protocol header
whose bit layout is dictated by a spec, bit fields are the right tool.
If your struct merely wants ``small flags'', prefer ordinary
\code{int} fields (or bitwise \code{\&} / \code{|} on a single
integer); the exact layout of bit fields across compilers is not
portable, and portability is a promise we keep in this course.
\end{ruleofthumb}

\section{Unions: Overlapping Storage}

A \code{union} is a struct whose members all occupy the \emph{same}
bytes. Its size is the size of its largest member, not the sum. Writing
one member overwrites the others, so a union is only as meaningful as
the convention that says which member is live. Pair a union with a
``kind'' tag and you have a \emph{tagged union}, the C ancestor of the
enum variant that powers trees, graphs, and protocol parsers everywhere.

That tag is an \code{enum}---short for \emph{enumeration}, a type that
names a set of integer constants. \code{enum kind \{ KIND\_INT,
KIND\_DOUBLE, KIND\_TEXT \};} declares three constants
(\code{KIND\_INT} is 0, the next 1, the next 2, unless you say
otherwise) and a type \code{kind} that holds one of them. It exists so
your code can say \code{v.tag = KIND\_INT} instead of the
incomprehensible \code{v.tag = 0}. In C an enum is just a set of named
integer constants, so the compiler does not police which values you
assign or compare---the names exist for humans, not for enforcement.
This is the C version of a named choice, and you will see it in every
parser and protocol handler later in the book.

\begin{verbatim}
   struct or union?  Same members, very different size.

   struct thing {                   union thing {
       int n;      // 4 bytes            int n;      // 4 bytes
       double d;   // 8 bytes            double d;   // 8 bytes
       char c[16]; // 16 bytes           char c[16]; // 16 bytes
   };               // 32 bytes       };              // 16 bytes

   A struct keeps every member:               A union reuses the space:
   +----+----+----+----+----+----+----+      +----+----+----+----+----+
   | n  | pad |      d      |    c[16]   |   |   n, or d, or c[16]    |
   +----+----+----+----+----+----+----+----+  +----+----+----+----+----+
   0    4   8                16     32        0                        16
\end{verbatim}

The standard trick is to record \emph{which} member is live in a
sibling field, and to treat writes to the union as a single operation.

\begin{lstlisting}[caption={A tagged union: kind says which member is live.}]
#include <stdio.h>

typedef enum kind {
    KIND_INT,
    KIND_DOUBLE,
    KIND_TEXT
} kind;

typedef struct value {
    kind tag;              /* which member of u is valid */
    union {
        int    as_int;
        double as_double;
        char   as_text[16];
    } u;
} value;

void value_print(const value *v) {
    /* We promise to read only the member named by v->tag. */
    if (v->tag == KIND_INT) {
        printf("%d\n", v->u.as_int);
    } else if (v->tag == KIND_DOUBLE) {
        printf("%f\n", v->u.as_double);
    } else {
        printf("%s\n", v->u.as_text);
    }
}

int main(void) {
    value v;

    v.tag = KIND_INT;
    v.u.as_int = 42;
    value_print(&v);

    v.tag = KIND_TEXT;
    v.u.as_text[0] = 'h';
    v.u.as_text[1] = 'i';
    v.u.as_text[2] = '\0';
    value_print(&v);

    return 0;
}
\end{lstlisting}

The danger is symmetrical and classic: read the wrong member and you
interpret one kind of bits as another, silently producing garbage
numbers or strings. The tag discipline is the price of memory sharing.
Used with discipline, tagged unions give us a single \code{value} type
for the key/value engines, expression trees, and configuration parsers
we build in later chapters.

\section{Practical Exercise: A Person Registry}

Let us assemble the chapter into a working registry: a linked list of
\code{struct person} records, built entirely with the techniques we
have seen. This is the skeleton on which the serious containers of
Chapter~4 will be built, so read its shape carefully: constructor,
push-front, traversal, destroy.

\begin{lstlisting}[caption={registry.c: a minimal linked list of persons.}]
#include <stdio.h>
#include <stdlib.h>
#include <string.h>

struct person {
    char name[32];
    int  age;
    struct person *next;      /* the link to the following person, or NULL */
};

/* Create a new person on the heap. Returns NULL on failure. */
struct person *person_new(const char *name, int age) {
    struct person *p;

    p = malloc(sizeof(struct person));
    if (p == NULL) {
        return NULL;
    }

    strncpy(p->name, name, sizeof(p->name) - 1);
    p->name[sizeof(p->name) - 1] = '\0';
    p->age = age;
    p->next = NULL;

    return p;
}

/* Insert p at the front of the list named by *head. */
void list_push_front(struct person **head, struct person *p) {
    p->next = *head;          /* p points at whatever was first before  */
    *head = p;                /* and *head now names p as the new first */
}

/* Print every person in the list. */
void list_print_all(const struct person *head) {
    const struct person *cur;

    cur = head;
    while (cur != NULL) {
        printf("%s, age %d\n", cur->name, cur->age);
        cur = cur->next;      /* advance to the next node */
    }
}

/* Free every node in the list. */
void list_free_all(struct person *head) {
    struct person *cur;
    struct person *next;

    cur = head;
    while (cur != NULL) {
        next = cur->next;     /* remember the successor before freeing  */
        free(cur);            /* free this node                          */
        cur = next;           /* move on                                */
    }
}

int main(void) {
    struct person *head = NULL;

    list_push_front(&head, person_new("ada", 36));
    list_push_front(&head, person_new("linus", 55));
    list_push_front(&head, person_new("ken", 81));

    list_print_all(head);

    list_free_all(head);
    head = NULL;   /* head is now dangling; clearing it is good hygiene */

    return 0;
}
\end{lstlisting}

Study \code{list\_push\_front} closely. It takes \code{struct person
**head} because it must \emph{rewrite} the caller's \code{head}
variable---the exact pointer-to-pointer pattern from Chapter~1.
And \code{list\_free\_all} saves \code{next} \emph{before} calling
\code{free}, because after \code{free(cur)} the node is gone and
reading \code{cur->next} would be a use-after-free. Every line here
exists to keep the arrows honest.

\section{Structures in Bytes}

The layouts you computed by hand earlier in this chapter, collected on
one page. On a 64-bit machine with 4-byte \code{int} and 8-byte
\code{double}:

\begin{verbatim}
   struct point { int x; int y; }        size  8   x:0  y:4
   struct entry (tag, count, level)      size 12   tag:0 count:4 level:8
   struct entry (count, level, tag)      size  8   count:0 level:4 tag:6
   struct rect  { point tl; point br; }  size 16   tl:0  br:8
   union thing  { int n; double d; char c[16]; }   size 16
   struct flags (three bit fields)       size  4   one storage unit
   struct message { uint len; char data[]; }       size  4  (FAM adds nothing)

   "field:off" is the byte offset a debugger reports; size is what
   sizeof() returns. Both come from the alignment rules, never from
   hand arithmetic.
\end{verbatim}

The three rules in byte terms: every field starts at a byte offset
divisible by its own alignment; the struct's size rounds up to a
multiple of its strictest member's alignment; and a union occupies
exactly its largest member. Bit fields are the exception to ``every
field has an address''---a 1-bit field has no byte address, only a bit
position inside its storage unit. And \code{offsetof(struct entry,
level)} being 8 while \code{sizeof(struct entry)} is 12 is the
padding: 4 bytes the compiler reserved for you. Those padding bytes
are exactly why Chapter~16 will refuse to \code{fwrite} a struct to
disk and will write a wire format instead.

\section{Summary}

We now possess the two great tools of C: addresses/pointers
(Chapter~1) and structured, contiguous records (this chapter). The
combination---structs that point at each other---is the raw material of
every data structure that follows.

\begin{verbatim}
   The toolbox so far:

   +--------------------------+-----------------------------------------+
   | idea                     | the tool                                |
   +--------------------------+-----------------------------------------+
   | an address               | pointer: int *p;                        |
   | a grouped record         | struct point { int x; int y; };         |
   | a short alias for a type | typedef struct point {...} point;       |
   | record on the heap       | malloc(sizeof(point));                  |
   | field through ptr        | p->x                                    |
   | read-only contract       | const point *p                          |
   | a struct inside a struct | point top_left;  (nested by value)      |
   | a chain of records       | struct node { ...; struct node *next; } |
   | header + payload         | char data[];   (flexible array member)  |
   | silent extra bytes       | alignment padding (order fields!)       |
   | one storage, many views  | union { ... }  + a tag                  |
   | rows or columns          | AoS vs SoA (choose by access pattern)   |
   +--------------------------+-----------------------------------------+
\end{verbatim}

In Chapter~4 we go to work. The Mega Chapter builds and dissects the
entire menagerie---linked lists in every flavor, skip lists, hash
tables, trees from $N$-ary to B+ trees, and graphs with edges and
weights---each one drawn with memory diagrams and implemented in this
explicit, pointer-honest style.
